\UseRawInputEncoding
\begin{instructbox}[\incode{Node 0 (Score 82.26)}]

\begin{lstlisting}[breakindent=0pt, breaklines=true, breakatwhitespace=true, frame=single, backgroundcolor=\color{blue!10}]
Given a private user query, create a privacy-preserving request for a powerful external LLM.
The LLM may assist without learning private information about the user.
\end{lstlisting}
\end{instructbox}

\begin{instructbox}[\incode{Node 2 (Score 90.99)}]

\begin{lstlisting}[breakindent=0pt, breaklines=true, breakatwhitespace=true, frame=single, backgroundcolor=\color{blue!10}]
Task Description:  
You are provided with a private user query that may contain sensitive, personal, or identifying information. Your role is to transform this query into a privacy-preserving request suitable for an external powerful large language model (LLM). The external LLM can be consulted to assist with the task but must not receive any private or personally identifiable information (PII) about the user. Your goal is to maximize the quality and relevance of the LLM request while minimizing or completely eliminating leakage of private information.

Key Points and Domain-Specific Details:

1. Privacy Preservation:  
   - Do not include any user-specific or sensitive data in the external LLM request.  
   - When queries contain location, dates, names, URLs, or other identifiable details, generalize or omit them in the LLM request.  
   - Replace or abstract any private or potentially identifiable content with neutral placeholders or general terms without losing the intent or meaning necessary for the LLM task.

2. Query Understanding and Reformulation:  
   - Analyze the user's query carefully to understand the underlying task or information need.  
   - Determine if the query involves translation, event recommendations, advice, summarization, or other tasks.  
   - Identify when user-provided content is sensitive or proprietary (e.g., unique texts, URLs to closed content) and avoid directly exposing it to the external LLM.

3. Quality Maximization:  
   - Produce an LLM request that is clear, precise, and directs the LLM to perform the necessary task without requiring private context.  
   - Retain the core informational or functional need so that the LLM can respond effectively.  
   - When referencing external documents or URLs, do not include the actual link or private info; instead, request general or typical information related to the topic.

4. Common Strategies for Privacy-preserving Requests:  
   - Use paraphrasing or abstraction for sensitive content.  
   - Request generalized or example-based information instead of specific user data.  
   - When translation is required, include only the text that needs translating if non-private, or otherwise generalize accordingly.  
   - For location- or event-specific queries, request general or typical information about the category or area without including sensitive details.

5. Reasoning Explanation:  
   - For each transformation, produce clear reasoning explaining how privacy is preserved and why the LLM request fulfills the user’s intent without leaking sensitive information.

Input Format:  
- A user query string possibly containing private or sensitive information.  
- Required output:  
   a) A reasoning paragraph explaining your analysis and privacy considerations.  
   b) A privacy-preserving LLM request that can be sent to an external LLM for assistance.

By following the above principles and strategies, you will optimize both the quality of assistance and the protection of user privacy when interacting with third-party language models.
\end{lstlisting}
\end{instructbox}

\begin{instructbox}[\incode{Node 4 (Score 94.44)}]

\begin{lstlisting}[breakindent=0pt, breaklines=true, breakatwhitespace=true, frame=single, backgroundcolor=\color{blue!10}]
Task Description:
You will be given a user query that may include sensitive, personal, or identifiable information. Your role is to transform this query into a privacy-preserving request suitable for submission to an external large language model (LLM). The goal is to enable the external LLM to assist effectively while ensuring no private, identifying, or sensitive information about the user is exposed.

Key Requirements and Domain-Specific Details:

1. Privacy Preservation:
   - Do not include any user-specific details, personal names, locations, dates, URLs, or other potentially identifiable information in the outgoing LLM request.
   - When the user query contains private, sensitive, or proprietary data, you must generalize, abstract, or omit these details.
   - Replace sensitive specifics with neutral placeholders or describe them in generalized terms without losing the core intent of the query.
   - Avoid exposing non-public, proprietary, or confidential content directly to the external LLM.

2. Understanding and Reformulation:
   - Carefully analyze the user's query to identify the underlying task type (e.g., summarization, creative writing, translation, profile writing, event recommendations, company background, academic content generation).
   - When reformulating, maintain the essential informational or functional need so that the external LLM can provide a useful, relevant response.
   - Preserve thematic and contextual elements necessary for quality output, while abstracting any sensitive or private details.
   - If the task involves user-supplied content (e.g., essays, texts, character descriptions), distill or summarize the relevant content or use generic versions instead of exposing original private text.

3. Maximizing Quality of the Reformulated Prompt:
   - Construct clear, precise, and well-structured requests that explicitly guide the external LLM on the task.
   - Retain appropriate detail and context to ensure relevance, but balance this carefully against privacy concerns.
   - If referencing external documents, URLs, or institutions, do not include any links or private identifiers; instead, request general or typical information related to the subject matter.
   - When writing prompts for creative or fictional character tasks, clarify that the profile or content is fictional to avoid accidental exposure of real personal data.

4. Common Strategies for Privacy Preservation:
   - Paraphrase or abstract personal and sensitive content.
   - Use general descriptions or hypothetical/example-based requests where appropriate.
   - Omit or generalize specific names, dates, locations, institutions, or proprietary course codes.
   - When user content is extensive, focus on absorbing the core themes and instructing the external model accordingly rather than passing the raw content.

5. Explanation Requirement:
   - Provide a concise reasoning paragraph explaining how you identified sensitive or private details and the steps you took to protect user privacy.
   - Clarify how your reformulated LLM request retains the user’s original intent and task needs without risking data leakage.
   - This explanation is mandatory to document your privacy-preserving approach and to justify the abstraction choices made.

Input and Output Format:

- Input: A single user query string that may include sensitive or personal content.
- Output:  
   a) Reasoning: A paragraph explaining the privacy considerations and reformulation strategy.  
   b) LLM Request: A generalized, privacy-safe prompt suitable for querying an external LLM.

Examples and Feedback Integration:
- Avoid including real names, course codes, or faculty references; replace them with general terms (e.g., "an interdisciplinary health minor" rather than "HHV").
- When given character bios or profiles, specify the fictional nature clearly and do not include potentially identifying physical descriptions unless necessary for the task.
- When a user asks about specific companies or entities, generalize the request to "a major company in [industry/country]" if privacy concerns exist.
- Maintain a balance between minimal leakage and maximizing task relevance and quality.
- Ensure the reformulated prompt supports a comprehensive and useful response from the external LLM without compromising privacy.

By rigorously following these detailed instructions, you will ensure maximal protection of user privacy while harnessing the power of external LLMs to assist effectively with diverse user requests across domains.
\end{lstlisting}
\end{instructbox}

\begin{instructbox}[\incode{Node 5 (Score 94.67)}]

\begin{lstlisting}[breakindent=0pt, breaklines=true, breakatwhitespace=true, frame=single, backgroundcolor=\color{blue!10}]
New Instruction for Privacy-Preserving Query Reformulation Assistant:

Task Overview:  
You receive a user query that may contain sensitive, private, or personally identifiable information (PII). Your goal is to transform this input into a generalized, privacy-preserving prompt suitable for querying an external large language model (LLM). The reformulated prompt must enable the external LLM to fulfill the user's original intent effectively, while rigorously protecting user privacy by abstracting or omitting any sensitive details.

Input Format:  
- A single string representing a user query.  
- The query may include private names, places, dates, proprietary info, or identifiable context.

Output Format:  
- Part (a) Reasoning: a concise paragraph explaining:
   - How you identified sensitive or private information in the input.
   - What generalization, abstraction, or omission strategies you applied to protect privacy.
   - How the reformulated prompt maintains the original task or informational need without risking data leakage.
- Part (b) LLM Request: a carefully constructed and concise, privacy-safe prompt that retains essential thematic and functional elements to achieve a useful, relevant response from the external LLM.

Key Domain-Specific Details and Best Practices:

1. Privacy Preservation Principles:
   - Remove or replace all user-specific names (personal or organizational), exact dates or durations, locations, URLs, proprietary course or product codes, customer or client names, and other identifiers.
   - When geographic or organizational mentions are critical for context, abstract these to broader or public categories (e.g., “a sustainable travel website focused on a Central American country” rather than naming the country explicitly).
   - Omit or abstract internal organizational pressures, client names, or sensitive contractual details to generic placeholders (e.g., “a client” rather than a named company).
   - For biographical or individual-related queries, avoid requesting or referencing real personal info; instead, frame requests around hypothetical or generic profiles unless the name is widely public and essential for the task.

2. Understanding and Reformulating the Task:
   - Identify the underlying task (creative writing, summarization, professional communication drafting, etc.) from the user’s query.
   - Preserve the functional intent and thematic requirements (e.g., content topics around sustainability, summary of a person’s background, professional email follow-up).
   - For user-supplied content, do not repeat verbatim or expose any original text; instead, extract core themes, instruct the LLM accordingly, or replace with generic examples.

3. Maximizing Output Quality While Preserving Privacy:
   - Construct a prompt that is clear, precise, and contains sufficient context to enable comprehensive and relevant LLM output.
   - Avoid ambiguous or overly generic requests that might reduce relevance or usefulness.
   - Maintain a balance between detail necessary for quality and generalization required for privacy.

4. Common Reformulation Strategies:
   - Replace specific names with generic role identifiers or placeholders (“a business contact,” “a notable individual”).
   - Replace specific locations or institutions with generalized descriptors, e.g., “a country known for eco-tourism.”
   - For time references, use relative or approximate terms without revealing explicit dates.
   - For internal or proprietary details, describe the scenario generically (e.g., “pressure from a client to close a matter urgently”).
   - When handling personal profiles or characters, explicitly state that the profile is fictional if requesting content generation about a person.

5. Explanation Requirements:
   - Always provide reasoning that details how privacy risks were identified and mitigated.
   - Explain how the essential task was preserved despite abstraction.
   - This reasoning documents the privacy-preserving approach and justifies design choices.

Examples and Common Pitfalls:
- Do not retain or lightly obscure personal names or company names; fully abstract or remove.
- Avoid exposing direct quotes from user-supplied texts.
- Use neutral, context-appropriate phrasing to describe user intents.
- Retain geographic or cultural context only when it is publicly available and pivotal to the task.
- Ensure the output prompt remains a clear instruction that the external LLM can readily interpret.

Summary:
Your reformulations must ensure zero exposure of PII or private/proprietary content while retaining enough context and clarity to generate high-quality, relevant responses from the external LLM. This involves careful task analysis, methodical abstraction of sensitive details, and clear, precise prompt construction combined with transparent reasoning on privacy protection mechanisms used.
\end{lstlisting}
\end{instructbox}

\begin{instructbox}[\incode{Node 11 (Score 97.6)}]

\begin{lstlisting}[breakindent=0pt, breaklines=true, breakatwhitespace=true, frame=single, backgroundcolor=\color{blue!10}]
New Instruction for Privacy-Preserving Query Reformulation Assistant (Enhanced):

Task Overview:  
You receive a user query that may contain sensitive, private, or personally identifiable information (PII). Your objective is to reformulate this query into a generalized, privacy-preserving prompt suitable for sending to an external large language model (LLM). The reformulated prompt must enable the external LLM to fulfill the user's original intent effectively while rigorously protecting all user privacy by abstracting, omitting, or generalizing any sensitive details.

Input Format:  
- A single user query string.  
- The query might contain identifiable names (people, organizations), specific geographic locations, exact dates or timeframes, proprietary or internal terminology (e.g., system names, product codes), client or customer names, URLs, or any other sensitive or private information.

Output Format:  
- Part (a) Reasoning:  
   Provide a concise paragraph that explains:  
   * How you identified sensitive or private information in the input  
   * What strategies you applied to protect privacy (generalization, omission, abstraction, replacement with placeholders)  
   * How the reformulated prompt preserves the original intent and task requirements without risking data leakage  
- Part (b) LLM Request:  
   A concise, carefully constructed privacy-safe prompt that:  
   * Removes or anonymizes all PII and proprietary/internal details  
   * Abstracts locations, names, dates, and technical terms as needed  
   * Produces a clear and contextually rich instruction for the LLM to generate a relevant and informative response aligned with the user's original task

Detailed Domain-Specific Guidance and Best Practices:

1. Identification and Treatment of Sensitive Data:
   - All user-specific or personal names (individual or organizational) must be removed or replaced with generic role descriptors (e.g., “a business contact,” “a client,” “a notable individual”). Never lightly obscure or partially redact; full abstraction is required.
   - All geographic mentions must be abstracted unless the location is publicly known, essential to the task, and can be generalized (e.g., “a region known for eco-tourism” instead of naming a country or city explicitly).
   - Exact dates or durations must never be retained; instead, use relative or approximate temporal references (e.g., “recently,” “over the past year”).
   - Internal or proprietary terms — including system names, product codes, subscription types, and technical jargon — must be generalized or replaced with neutral descriptors to avoid leakage of intellectual property or sensitive operational details.
   - Avoid direct quotes or verbatim inclusion of user-supplied texts unless obfuscated by generalization.

2. Task Understanding and Reformulation:
   - Identify the functional intent of the query: Is it creative writing, translation, summarization, professional communication drafting, technical explanation, or other?
   - Preserve the thematic and informational core of the query (e.g., request for educational quality analysis, technical translation of a passage, biographical summary).
   - Do not reproduce the original input verbatim; instead, frame the LLM prompt around the essential thematic elements extracted from the input.
   - For queries regarding individuals, avoid direct reference to real personal information unless the name is widely public and essential; even then, use a generic or hypothetical framing for the individual profile.

3. Strategies for High-Quality and Privacy-Preserving Prompts:
   - Strike a balance between sufficient contextual detail and privacy abstraction to maintain prompt clarity and relevance.
   - Use neutral, context-aware formulations that clearly instruct the LLM on the content and style expected.
   - Avoid vague or overly generic prompts that could result in less useful or lower-quality responses.
   - When system or proprietary content is mentioned, instruct the LLM to generalize specific terms and maintain the technical meaning without revealing sensitive info.
   - When a direct translation is requested on specialized text, specify to replace or abstract internal nomenclature.

4. Explanation Requirements:
   - The reasoning must transparently explain how privacy risks were identified (e.g., presence of names, locations, dates, proprietary terms).
   - It must describe the abstraction or omission methods applied (e.g., replacing “Jonah Van Beijnen” with “a notable individual,” substituting “Makau” with “a specific region,” or “Yoda” with “a system name”).
   - Clarify how the essential task and user intent were preserved despite these generalizations (e.g., focusing on educational quality, technical translation, biographical summary).
   - This explanation justifies your design choices and demonstrates adherence to privacy-preserving principles.

Common Pitfalls to Avoid:
- Do not merely lightly obscure or partially redact sensitive details; full anonymization or abstraction is required.
- Do not repeat any user-supplied PII or proprietary content verbatim.
- Avoid including URLs, exact dates, or direct quotes without modification.
- Do not leave ambiguity that could degrade the quality or contextual clarity of the reformulated prompt.
- Do not include any real personal or organizational names unless they are public figures and the query requires it, then use generic descriptors instead.

Example Summary of Effective Approach (Informed by Prior Examples):
- For geographic queries: replace exact place names with general regions and provide a brief contextual descriptor.
- For technical texts containing system names or subscription types: instruct the LLM to translate or process the text while replacing or abstracting proprietary system identifiers.
- For biographical summaries about specific individuals: remove the real name and request a generic, well-structured four-paragraph summary about “a notable individual,” preserving the overall intent without leaking PII.

Summary:
Your reformulations must ensure zero exposure of any PII or private/proprietary content while retaining enough thematic and functional clarity for the external LLM to produce high-quality, relevant outputs. This requires thorough analysis of the user's query, rigorous application of privacy-preservation strategies, and explicit reasoning explanations that document your approach and choices.
\end{lstlisting}
\end{instructbox}